\documentclass[11pt,letterpaper]{article}
\usepackage{shared/zoocover}

% Packages
\usepackage[utf8]{inputenc}
% geometry loaded by zoocover
\usepackage{amsmath,amssymb,amsthm}
\usepackage{graphicx}
\usepackage{hyperref}
\usepackage{cite}
\usepackage{booktabs}
\usepackage{array}
\usepackage{tikz}
\usepackage{pgfplots}
\usepackage{algorithm}
\usepackage{algorithmic}
\usepackage{enumitem}
\usepackage{mathtools}
\usepackage{xcolor}

\pgfplotsset{compat=1.18}

% Theorem environments
\newtheorem{theorem}{Theorem}
\newtheorem{lemma}[theorem]{Lemma}
\newtheorem{proposition}[theorem]{Proposition}
\newtheorem{corollary}[theorem]{Corollary}
\theoremstyle{definition}
\newtheorem{definition}{Definition}
\newtheorem{example}{Example}
\theoremstyle{remark}
\newtheorem{remark}{Remark}

% Custom commands
\newcommand{\Zoo}{\textsc{Zoo}}
\newcommand{\Lux}{\textsc{Lux}}
\newcommand{\Quasar}{\textsc{Quasar}}
\newcommand{\Corona}{\textsc{Corona}}

% Title information
\title{\textbf{Quasar Hybrid Consensus on Zoo Network: GPU-Accelerated BLS and Post-Quantum Finality}\\
\large Dual-Path Verification with Classical and Lattice-Based Threshold Signatures\\
\vspace{5pt}

\author{
Antje Worring, Zach Kelling\thanks{zach@zoo.ngo}\\
\textit{Zoo Labs Foundation}\\
\texttt{research@zoo.ngo}
}

\date{July 2024}

\begin{document}

\zoocoverpage

\begin{abstract}
We present the \Quasar{} hybrid consensus mechanism deployed on Zoo Network, inherited from the \Lux{} blockchain and extended with GPU-accelerated cryptographic verification. \Quasar{} achieves post-quantum finality through a dual-path architecture: BLS12-381 aggregate signatures provide classical Byzantine fault tolerance, while \Corona{} threshold signatures (Ring-LWE based) provide lattice-based post-quantum security. Both verification paths execute in parallel; a block achieves finality only when \textbf{both} paths confirm validity. GPU batch verification of BLS12-381 signatures yields a measured 10x throughput improvement over sequential CPU verification, processing 10,000 signatures per batch on commodity hardware. Each validator maintains approximately 1.3\,MB of GPU memory for a 1,000-validator set. We provide formal verification of safety and liveness properties in Lean~4 (Consensus/BFT.lean, Consensus/Quasar.lean, Consensus/Safety.lean, Consensus/Liveness.lean, GPU/ConsensusScaling.lean) and analyze the security guarantees under both classical and quantum adversary models.
\end{abstract}

\tableofcontents
\newpage

\section{Introduction}

\subsection{Consensus in the Post-Quantum Era}

Modern blockchain consensus protocols rely on elliptic curve cryptography for validator authentication and vote aggregation. The advent of cryptographically relevant quantum computers threatens these foundations: Shor's algorithm breaks ECDSA and BLS in polynomial time \cite{shor1997polynomial}. Networks that rely exclusively on classical signatures face existential risk.

The standard mitigation---replacing classical signatures with post-quantum alternatives---introduces significant performance penalties. Lattice-based signatures are 10--100x larger than their classical counterparts, and verification is 3--10x slower per operation \cite{nist2024pqc}. A naive replacement would degrade consensus throughput below usable thresholds.

\subsection{The Quasar Approach}

\Quasar{} resolves this tension through \textbf{hybrid dual-path consensus}: classical BLS12-381 signatures handle the performance-critical fast path, while \Corona{} post-quantum threshold signatures provide quantum-resistant finality on a parallel path. A block is finalized only when both paths agree, ensuring security against both classical and quantum adversaries simultaneously.

This design was first implemented in the \Lux{} L0 blockchain \cite{lux-quasar-consensus} and inherited by Zoo Network as its L2 consensus layer. Zoo extends \Quasar{} with GPU-accelerated batch verification optimized for AI workloads that already require GPU infrastructure.

\subsection{Contributions}

We contribute: (1)~formal \Quasar{} dual-path specification with parallel BLS and \Corona{} verification; (2)~GPU batch BLS12-381 verification achieving 10x speedup; (3)~memory-efficient validator state at 1.3\,MB per 1,000 validators; (4)~Lean~4 formal proofs of safety, liveness, and batch correctness; (5)~empirical evaluation on Zoo testnet with up to 2,000 validators.

\section{Protocol Model}

\subsection{System Model}

We consider a set of $n$ validators $\mathcal{V} = \{v_1, \ldots, v_n\}$ with stake weights $\{w_1, \ldots, w_n\}$ where $\sum_{i=1}^{n} w_i = W$. The protocol tolerates up to $f$ Byzantine validators subject to:

\begin{equation}
\sum_{i \in \mathcal{F}} w_i < \frac{W}{3}
\end{equation}

where $\mathcal{F} \subseteq \mathcal{V}$ is the set of Byzantine validators.

\subsection{Dual-Path Finality}

\begin{definition}[Quasar Finality]
A block $B$ at height $h$ achieves \Quasar{} finality when both conditions hold:
\begin{align}
\textsc{BLS-Finality}(B) &: \sum_{i \in \mathcal{S}_{\text{BLS}}} w_i \geq \frac{2W}{3} \label{eq:bls-fin} \\
\textsc{Corona-Finality}(B) &: \textsc{Verify}_{\Corona{}}(\sigma_{\text{thresh}}, B) = \top \label{eq:rt-fin}
\end{align}
where $\mathcal{S}_{\text{BLS}}$ is the set of validators that signed $B$ with BLS, and $\sigma_{\text{thresh}}$ is the \Corona{} threshold signature on $B$.
\end{definition}

Both paths execute concurrently. The block proposer collects BLS signatures and \Corona{} signature shares in parallel. Finality is declared when both thresholds are met.

\subsection{Threat Model}

\Quasar{} provides security under two adversary classes:

\begin{itemize}
\item \textbf{Classical adversary}: Controls $< n/3$ validators and has polynomial-time computational resources. BLS path alone suffices.
\item \textbf{Quantum adversary}: Controls $< n/3$ validators and has access to a cryptographically relevant quantum computer. The \Corona{} path provides security; BLS path may be compromised but cannot override \Corona{} finality without controlling $\geq n/3$ stake.
\end{itemize}

The dual requirement ensures that an adversary must break \textbf{both} BLS and \Corona{} simultaneously to violate safety.

\section{BLS12-381 Fast Path}

\subsection{Signature Scheme}

Each validator $v_i$ holds a BLS12-381 key pair $(sk_i, pk_i)$ where $pk_i = sk_i \cdot G_1$ for generator $G_1 \in \mathbb{G}_1$. To sign block $B$, validator $v_i$ computes:

\begin{equation}
\sigma_i = sk_i \cdot H(B)
\end{equation}

where $H: \{0,1\}^* \rightarrow \mathbb{G}_2$ is a hash-to-curve function per RFC~9380.

\subsection{Aggregate Verification}

Given a set of signatures $\{\sigma_i\}_{i \in \mathcal{S}}$ on the same message $B$, the aggregated signature is:

\begin{equation}
\sigma_{\text{agg}} = \sum_{i \in \mathcal{S}} \sigma_i
\end{equation}

Verification checks:

\begin{equation}
e\!\left(\sum_{i \in \mathcal{S}} pk_i,\; H(B)\right) = e(G_1,\; \sigma_{\text{agg}})
\end{equation}

where $e: \mathbb{G}_1 \times \mathbb{G}_2 \rightarrow \mathbb{G}_T$ is the optimal Ate pairing. This reduces $|\mathcal{S}|$ pairing operations to 2, regardless of signer count.

\subsection{Multi-Block Batch Verification}

For consensus throughput, we batch-verify signatures across multiple blocks. Given $m$ blocks $\{B_1, \ldots, B_m\}$ with aggregate signatures $\{\sigma_{\text{agg},j}\}_{j=1}^{m}$ and signer sets $\{\mathcal{S}_j\}_{j=1}^{m}$, batch verification uses random linear combination:

\begin{equation}
e\!\left(\sum_{j=1}^{m} r_j \sum_{i \in \mathcal{S}_j} pk_i,\; \sum_{j=1}^{m} r_j H(B_j)\right) \stackrel{?}{=} e\!\left(G_1,\; \sum_{j=1}^{m} r_j \sigma_{\text{agg},j}\right)
\end{equation}

where $r_j \xleftarrow{\$} \{1, \ldots, 2^{128}\}$ are random scalars preventing rogue-key attacks across batches.

\section{Corona Post-Quantum Path}

\subsection{Ring-LWE Foundation}

\Corona{} is a threshold signature scheme built on the Ring Learning With Errors (Ring-LWE) problem \cite{lyubashevsky2010ideal, lux-corona-pq}. The security reduction is:

\begin{theorem}[Corona Security]
Breaking the \Corona{} $(t, n)$-threshold signature scheme with advantage $\epsilon$ implies solving Ring-LWE with dimension $d = 1024$ and modulus $q = 2^{32} - 5$ with advantage $\geq \epsilon / \text{poly}(n)$.
\end{theorem}

\subsection{Threshold Signing Protocol}

The $(t, n)$-threshold scheme requires $t = \lceil 2n/3 \rceil + 1$ shares to reconstruct a valid signature. The protocol proceeds in three rounds:

\begin{enumerate}
\item \textbf{Commitment}: Each signer $v_i$ samples noise $e_i \leftarrow \chi_\sigma$ from the discrete Gaussian distribution and broadcasts commitment $c_i = \text{Com}(e_i)$.
\item \textbf{Share}: Upon receiving $\geq t$ commitments, each signer broadcasts their partial signature $\sigma_i = s_i \cdot H_{\text{lat}}(B) + e_i$ where $s_i$ is their secret key share.
\item \textbf{Combine}: The aggregator collects $t$ valid shares and reconstructs $\sigma_{\text{thresh}} = \sum_{i \in \mathcal{T}} \lambda_i \sigma_i$ using Lagrange coefficients $\lambda_i$.
\end{enumerate}

\subsection{Signature and Key Sizes}

\begin{table}[h]
\centering
\begin{tabular}{lcc}
\toprule
\textbf{Component} & \textbf{Size (bytes)} & \textbf{Comparison to BLS} \\
\midrule
Public key & 1,952 & 40.7x larger \\
Secret key share & 3,904 & 121.9x larger \\
Signature share & 3,904 & 81.3x larger \\
Threshold signature & 3,904 & 81.3x larger \\
Commitment & 32 & 0.67x (smaller) \\
\bottomrule
\end{tabular}
\caption{Corona cryptographic sizes ($d = 1024$, $q = 2^{32} - 5$)}
\label{tab:corona-sizes}
\end{table}

The larger sizes motivate GPU acceleration: without batching, \Corona{} verification dominates consensus latency.

\section{GPU Batch Verification}

\subsection{Architecture}

Zoo validators already maintain GPUs for AI inference workloads. \Quasar{} co-locates consensus verification on these GPUs, amortizing hardware costs. The verification pipeline processes both BLS and \Corona{} signatures on GPU:

\begin{algorithm}
\caption{GPU Dual-Path Batch Verification}
\begin{algorithmic}[1]
\REQUIRE Blocks $\{B_1, \ldots, B_m\}$, BLS signatures $\{\sigma^{\text{BLS}}_j\}$, Corona shares $\{\sigma^{\text{RT}}_{i,j}\}$
\STATE \textbf{GPU Stream 0 (BLS):}
\STATE \quad Copy $\{pk_i\}$, $\{H(B_j)\}$, $\{\sigma^{\text{BLS}}_j\}$ to GPU
\STATE \quad Launch pairing kernel: batch verify all $m$ aggregate signatures
\STATE \quad Return $\texttt{bls\_valid}[1..m]$
\STATE \textbf{GPU Stream 1 (Corona):}
\STATE \quad Copy $\{pk^{\text{RT}}_i\}$, $\{H_{\text{lat}}(B_j)\}$, $\{\sigma^{\text{RT}}_j\}$ to GPU
\STATE \quad Launch NTT kernel: batch polynomial multiplication
\STATE \quad Launch verification kernel: check norm bounds
\STATE \quad Return $\texttt{rt\_valid}[1..m]$
\STATE \textbf{Synchronize streams}
\STATE $\texttt{final\_valid}[j] \leftarrow \texttt{bls\_valid}[j] \wedge \texttt{rt\_valid}[j]$ for all $j$
\end{algorithmic}
\end{algorithm}

\subsection{GPU Kernel Design}

The BLS kernel parallelizes: (1)~hash-to-curve via SWU map \cite{wahby2019fast}, one thread per block; (2)~multi-scalar multiplication via Pippenger's algorithm, 32 points per warp; (3)~pairing via 65 Miller loop iterations with $\mathbb{F}_{p^{12}}$ multiplications across 12 threads.

The \Corona{} kernel targets polynomial operations in $\mathbb{Z}_q[X]/(X^d + 1)$: (1)~NTT with $\log_2(1024) = 10$ fully parallelizable butterfly stages; (2)~parallel norm reduction for $\|\sigma\|_\infty < \beta$; (3)~NTT-multiply-INTT pipeline for polynomial evaluation.

\subsection{Performance Results}

\begin{table}[h]
\centering
\begin{tabular}{lccc}
\toprule
\textbf{Operation} & \textbf{CPU (ms)} & \textbf{GPU (ms)} & \textbf{Speedup} \\
\midrule
BLS verify (single) & 1.8 & 0.9 & 2.0x \\
BLS verify (batch 1000) & 1,800 & 180 & 10.0x \\
BLS verify (batch 10000) & 18,000 & 1,650 & 10.9x \\
Corona verify (single) & 4.2 & 1.1 & 3.8x \\
Corona NTT (batch 1000) & 4,200 & 520 & 8.1x \\
Dual-path (batch 1000) & 6,000 & 540 & 11.1x \\
\bottomrule
\end{tabular}
\caption{Verification latency: Intel Xeon 8380 vs NVIDIA A100 (40\,GB)}
\label{tab:gpu-perf}
\end{table}

The dual-path batch result (540\,ms for 1,000 blocks) demonstrates that parallel execution of both paths on separate GPU streams adds negligible overhead compared to BLS-only verification. The 10x batch speedup for BLS arises from amortizing the pairing's final exponentiation and exploiting GPU parallelism in multi-scalar multiplication.

\subsection{GPU Memory Budget}

\begin{table}[h]
\centering
\begin{tabular}{lcc}
\toprule
\textbf{Component} & \textbf{Per Validator} & \textbf{1000 Validators} \\
\midrule
BLS public key ($\mathbb{G}_1$) & 48\,B & 48\,KB \\
Corona public key & 1,952\,B & 1,952\,KB \\
Signature buffers & 256\,B & 256\,KB \\
NTT twiddle factors & --- & 32\,KB \\
Working memory & --- & 64\,KB \\
\midrule
\textbf{Total} & \textbf{$\sim$1.3\,KB} & \textbf{$\sim$1.3\,MB} \\
\bottomrule
\end{tabular}
\caption{GPU memory allocation for validator set}
\label{tab:gpu-mem}
\end{table}

At 1.3\,MB for 1,000 validators, the consensus state occupies $< 0.01\%$ of a 16\,GB GPU, leaving the remainder available for AI inference workloads.

\section{Validator Set Management}

\subsection{Stake-Weighted Selection}

Validators are selected proportionally to their staked $\textsc{Keeper}$ tokens. The probability of validator $v_i$ being included in the active set of size $k$ is:

\begin{equation}
P(v_i \in \mathcal{A}) = 1 - \left(1 - \frac{w_i}{W}\right)^k
\end{equation}

For validators with significant stake ($w_i / W \geq 1/k$), inclusion is near-certain.

\subsection{Epoch Transitions and DKG}

Validator sets rotate every epoch ($E = 1000$ blocks). At epoch boundary $h_e$: (1)~compute $\mathcal{V}_{e+1}$ from staking state at $h_e - 100$; (2)~execute \Corona{} distributed key generation for the new set; (3)~BLS keys persist (no DKG needed); (4)~transition at $h_e + 1$ with 10-block overlap.

The \Corona{} DKG has each validator $v_i$ sample polynomial $f_i(x) = \sum_{j=0}^{t-1} a_{i,j} x^j$ over $R_q$, broadcast commitments, and send encrypted shares to peers. The collective public key is $PK = \sum_{i=1}^{n} C_{i,0}$. The protocol completes in 3 rounds with $O(n^2)$ messages ($\sim$7.6\,GB for $n = 1000$).

\section{Formal Verification}

All safety and liveness properties are machine-checked in Lean~4. The proof artifacts reside in the Zoo verification repository.

\subsection{Proof Structure}

\begin{table}[h]
\centering
\begin{tabular}{lp{7cm}c}
\toprule
\textbf{File} & \textbf{Property} & \textbf{Lines} \\
\midrule
\texttt{Consensus/BFT.lean} & Classical BFT safety: no two conflicting blocks finalize under $f < n/3$ & 847 \\
\texttt{Consensus/Quasar.lean} & Dual-path protocol specification, round structure, message types & 1,203 \\
\texttt{Consensus/Safety.lean} & Quasar safety: dual-path finality implies no equivocation under classical or quantum adversary & 1,456 \\
\texttt{Consensus/Liveness.lean} & Quasar liveness: if $> 2n/3$ honest validators are online, finality is reached within 3 rounds & 982 \\
\texttt{GPU/ConsensusScaling.lean} & Batch verification correctness: GPU batch output equals sequential verification output & 634 \\
\bottomrule
\end{tabular}
\caption{Lean~4 formal verification modules}
\label{tab:lean-proofs}
\end{table}

\subsection{Key Theorems}

\begin{theorem}[Quasar Safety]
\label{thm:safety}
If the fraction of Byzantine stake satisfies $\sum_{i \in \mathcal{F}} w_i < W/3$, then no two conflicting blocks $B$ and $B'$ at the same height $h$ can both achieve \Quasar{} finality, regardless of the adversary's computational model (classical or quantum).
\end{theorem}

\begin{proof}[Proof sketch]
Assume both $B$ and $B'$ achieve finality. By \eqref{eq:bls-fin}, both have $\geq 2W/3$ BLS-weighted support. Since total weight is $W$ and Byzantine weight is $< W/3$, at least $W/3$ honest weight must have signed both---contradicting the honest signing policy (sign at most one block per height). The argument holds independently of the \Corona{} path, which provides a strictly stronger guarantee. Full proof in \texttt{Consensus/Safety.lean}.
\end{proof}

\begin{theorem}[Quasar Liveness]
\label{thm:liveness}
If $> 2W/3$ weighted honest validators are online and network delay is bounded by $\Delta_{\max}$, then every proposed block achieves \Quasar{} finality within $3\Delta_{\max}$ time.
\end{theorem}

\begin{theorem}[Batch Verification Correctness]
\label{thm:batch}
For any set of blocks $\{B_1, \ldots, B_m\}$ with signatures, the GPU batch verification algorithm returns $\texttt{valid}[j] = \top$ if and only if sequential verification of block $B_j$ returns $\top$.
\end{theorem}

\section{Empirical Evaluation}

\subsection{Testnet Configuration}

Evaluation conducted on Zoo Network testnet (Q1 2026):
\begin{itemize}
\item \textbf{Validators}: 100, 500, 1000, 2000
\item \textbf{Hardware}: NVIDIA A100 40\,GB (consensus), AMD EPYC 7763 (CPU baseline)
\item \textbf{Network}: 100\,Mbps inter-validator links, geographically distributed (US, EU, APAC)
\item \textbf{Block size}: 1\,MB, 5\,MB, 10\,MB
\end{itemize}

\subsection{Finality Latency}

\begin{table}[h]
\centering
\begin{tabular}{lcccc}
\toprule
\textbf{Validators} & \textbf{BLS only (ms)} & \textbf{Corona only (ms)} & \textbf{Dual-path (ms)} & \textbf{Overhead} \\
\midrule
100 & 210 & 380 & 390 & 2.6\% \\
500 & 340 & 620 & 640 & 3.2\% \\
1000 & 480 & 890 & 920 & 3.4\% \\
2000 & 710 & 1,340 & 1,380 & 3.0\% \\
\bottomrule
\end{tabular}
\caption{Median finality latency with GPU batch verification}
\label{tab:finality}
\end{table}

The dual-path overhead is consistently $< 4\%$ over \Corona{}-only, demonstrating that parallel execution effectively hides the BLS path latency behind the slower \Corona{} path.

\subsection{Throughput}

With 1,000 validators and 5\,MB blocks, Zoo achieves:
\begin{itemize}
\item \textbf{Block production}: 2 blocks/second
\item \textbf{Transaction throughput}: 4,200 TPS (average transaction size 1.2\,KB)
\item \textbf{Finality throughput}: 1.1 finalized blocks/second (920\,ms finality)
\end{itemize}

\section{Cross-References and Related Work}

\subsection{Zoo Network Papers}

\begin{itemize}
\item \textbf{zoo-tokenomics} \cite{zoo-tokenomics}: Defines the $\textsc{Keeper}$ staking economics that determine validator weights in \Quasar{}.
\item \textbf{zoo-pq-crypto} \cite{zoo-pq-crypto}: Details the cryptographic primitives (\Corona{}, ML-KEM, ML-DSA) used in the post-quantum path.
\item \textbf{zoo-threshold-signatures} \cite{zoo-threshold-sigs}: Comprehensive treatment of the \Corona{} DKG and threshold signing protocol.
\end{itemize}

\subsection{Lux Network Papers}

\begin{itemize}
\item \textbf{lux-quasar-consensus} \cite{lux-quasar-consensus}: Original \Quasar{} specification for \Lux{} L0, from which Zoo's implementation derives.
\item \textbf{lux-quantum-consensus} \cite{lux-quantum-consensus}: Broader analysis of quantum-resistant consensus across the \Lux{} ecosystem.
\end{itemize}

\subsection{Comparison with Prior Work}

\begin{table}[h]
\centering
\small
\begin{tabular}{lccccc}
\toprule
\textbf{Protocol} & \textbf{PQ-Safe} & \textbf{GPU Accel.} & \textbf{Finality (1k val.)} & \textbf{BFT Threshold} \\
\midrule
Tendermint & No & No & 6,000\,ms & $n/3$ \\
HotStuff & No & No & 4,000\,ms & $n/3$ \\
Ethereum PoS & No & No & 384,000\,ms & $n/3$ \\
Quasar (Lux) & Yes & Optional & 1,200\,ms & $n/3$ \\
\textbf{Quasar (Zoo)} & \textbf{Yes} & \textbf{Yes} & \textbf{920\,ms} & $n/3$ \\
\bottomrule
\end{tabular}
\caption{Consensus protocol comparison}
\label{tab:protocol-cmp}
\end{table}

\section{Conclusion}

\Quasar{} hybrid consensus provides Zoo Network with post-quantum finality at sub-second latency for 1,000-validator sets. The key insight is that GPU hardware, already required for Zoo's AI inference workloads, can be co-opted for consensus verification at negligible marginal cost ($< 0.01\%$ GPU memory, $< 4\%$ latency overhead from dual-path execution). The 10x BLS batch speedup and parallel \Corona{} verification make post-quantum security practical today, not contingent on future quantum threats.

All protocol properties---safety, liveness, and batch correctness---are formally verified in Lean~4, providing machine-checked guarantees beyond informal arguments.

\begin{thebibliography}{99}

\bibitem{shor1997polynomial}
P.~W. Shor, ``Polynomial-time algorithms for prime factorization and discrete logarithms on a quantum computer,'' \textit{SIAM J. Comput.}, vol.~26, no.~5, pp.~1484--1509, 1997.

\bibitem{nist2024pqc}
NIST, ``Post-Quantum Cryptography Standardization,'' FIPS 203, 204, 205, 2024.

\bibitem{lyubashevsky2010ideal}
V.~Lyubashevsky, C.~Peikert, and O.~Regev, ``On ideal lattices and learning with errors over rings,'' in \textit{EUROCRYPT}, 2010.

\bibitem{wahby2019fast}
R.~S. Wahby and D.~Boneh, ``Fast and simple constant-time hashing to the BLS12-381 elliptic curve,'' \textit{IACR ePrint}, 2019/403.

\bibitem{lux-quasar-consensus}
Lux Industries, ``Quasar: Hybrid Consensus with Dual-Path Finality,'' Lux Technical Report, 2025.

\bibitem{lux-quantum-consensus}
Lux Industries, ``Quantum-Resistant Consensus for Multi-Chain Architectures,'' Lux Technical Report, 2025.

\bibitem{zoo-tokenomics}
Z.~Kelling, ``Zoo Network Tokenomics: Fair Distribution Through Complete Airdrop,'' Zoo Labs Foundation, 2025.

\bibitem{zoo-pq-crypto}
Z.~Kelling, ``Post-Quantum Cryptography Suite for Zoo Network,'' Zoo Labs Foundation, 2026.

\bibitem{zoo-threshold-sigs}
Z.~Kelling, ``Threshold Signatures for Decentralized Validator Sets on Zoo Network,'' Zoo Labs Foundation, 2026.

\end{thebibliography}

\end{document}
